\documentclass{article}

\input{../definitions.tex}

\usepackage{mathpartir}
\newdef{problem}
\usepackage[a4paper,margin=1cm,landscape]{geometry}
\newpf{solution}{}
%\usepackage[a4paper]{geometry}
%\NewDocumentEnvironment{solution}{+b}{}{}

\ExplSyntaxOn
\NewDocumentEnvironment{problemdef}{+b}
{\tl_gset:Nn \g_tmpa_tl {#1}}
{
  \begin{problem}
    \begin{mathpar}
      \tl_use:N \g_tmpa_tl
    \end{mathpar}
  \end{problem}
  \begin{solution}
    %\let\emptybox{\boxed{##1}}
    \DeclareDocumentCommand{\emptybox}{m}{\boxed{##1}}
    \begin{mathpar}
      \tl_use:N \g_tmpa_tl
    \end{mathpar}
  \end{solution}
}
\ExplSyntaxOff

\newcommand{\seq}[3]{#1 ⊢ #2 : #3}
\def\emp{\bullet}
\newcommand{\lam}[3]{λ #1 : #2.\: #3}
\renewcommand{\pair}[2]{\p{#1, #2}}
\newcommand{\match}[5]{\textsf{match}\:{#1}\:| \inl {#2} ↦ {#3} | \inr {#4} ↦ {#5}}
\newcommand{\emptybox}[1]{\boxed{\phantom{#1}}}
\newcommand{\nullaryrule}[2]{
  \ExpandArgs{c}\DeclareDocumentCommand{#1}{m}{\inferrule*[right=\rlap{\(#2\)}] { } {##1}}
}
\newcommand{\unaryrule}[2]{
  \ExpandArgs{c}\DeclareDocumentCommand{#1}{mm}{\inferrule*[right=\rlap{\(#2\)}] {##2} {##1}}
}
\newcommand{\binaryrule}[2]{
  \ExpandArgs{c}\DeclareDocumentCommand{#1}{mmm}{\inferrule*[right=\rlap{\(#2\)}] {##2 \\ ##3} {##1}}
}
\newcommand{\ternaryrule}[2]{
  \ExpandArgs{c}\DeclareDocumentCommand{#1}{mmmm}{\inferrule*[right=\rlap{\(#2\)}] {##2 \\ ##3 \\ ##4} {##1}}
}
%\newcommand{\ax}[1]{\inferrule*[right=\textsc{Ax}] { } {#1}}
\nullaryrule{ax}{\textsc{Ax}}
\binaryrule{prodI}{×\textsc{I}}
\unaryrule{prodEL}{×\textsc{E}₁}
\unaryrule{prodER}{×\textsc{E}₂}
\unaryrule{funI}{→\textsc{I}}
\binaryrule{funE}{→\textsc{E}}
\unaryrule{sumIL}{+\textsc{I}{}₁}
\unaryrule{sumIR}{+\textsc{I}{}₂}
\ternaryrule{sumE}{+\textsc{E}}


\newcommand{\pr}[2]{\textsf{pr}_{#1} {#2}}

\begin{document}

\section{Products and function types}

\subsection{Problems}


\begin{problemdef}
  \funI
  {\seq \emp {\lam x A {\lam y B {\pair x y}}} {A → B → A×B}}
  {\funI
    {\seq {x : A} {{\lam y B {\pair x y}}} {\emptybox {B → A×B}}}
    {\prodI
      {\seq {\emptybox {x : A, y : B}} {\emptybox {\pair x y}} {A×B}}
      {\ax {\seq {x : A, y : B} x A}}
      {\ax {\emptybox {\seq {x : A, y : B} y B}}}
    }
  }
\end{problemdef}

\begin{problem}
  Construct a term of type \(A×B → B×A\) in the empty context.
\end{problem}
\begin{solution}
  \begin{mathpar}
    \funI
    {\seq \emp {\lam x {A×B} {\pair {\pr2 x} {\pr1 x}}} {A×B → B×A}}
    {\prodI
      {\seq {x : A×B} {\pair {\pr2 x} {\pr1 x}} {B×A}}
      {\prodER
        {\seq {x : A×B} {\pr2 x} B}
        {\ax {\seq {x : A×B} {x} {A×B}}}
      }
      {\prodEL
        {\seq {x : A×B} {\pr1 x} A}
        {\ax {\seq {x : A×B} {x} {A×B}}}
      }
    }
  \end{mathpar}
\end{solution}

\begin{problem}
  Construct a term of type \(\p{A → B → C} → B → A → C\) in the empty context.
\end{problem}
\begin{solution}
  \begin{mathpar}
    \funI
    {\seq \emp {\lam f {A → B → C} {\lam y B {\lam x A {f x y}}}} {\p{A → B → C} → B → A → C}}
    {\funI
      {\seq {f : A → B → C} {\lam y B {\lam x A {f x y}}} {B → A → C}}
      {\funI
        {\seq {f : A → B → C, y : B} {\lam x A {f x y}} {A → C}}
        {\funE
          {\seq {Γ ≔ f : A → B → C, y : B, x : A} {f x y} {C}}
          {\funE
            {\seq Γ {f x} {B → C}}
            {\ax {\seq Γ f {A → B → C}}}
            {\ax {\seq Γ x A}}
          }
          {\ax {\seq Γ y B}}
        }
      }
    }
  \end{mathpar}
\end{solution}

\begin{problem}
  Construct a term of type \(A+B → B+A\) in the empty context.
\end{problem}
\begin{solution}
  \begin{mathpar}
    \funI
    {\seq \emp {\lam t {A+B} {\match t x {\inr x} y {\inl y}}} {A+B → B+A}}
    {\sumE
      {\seq {t : A+B} {\match t x {\inr x} y {\inl y}} {B+A}}
      {\ax {\seq {t : A+B} t A+B}}
      {\sumIR
        {\seq {t : A+B, x : A} {\inr x} {B+A}}
        {\ax {\seq {t : A+B, x : A} x {A}}}
      }
      {\sumIL
        {\seq {t : A+B, y : B} {\inl y} {B+A}}
        {\ax {\seq {t : A+B, y : B} y B}}
      }
    }
  \end{mathpar}
\end{solution}


\begin{problem}
  Construct a term of type \(A×(B+C) → A×B + A×C\) in the empty context.
\end{problem}
\begin{solution}
  \begin{mathpar}
    \funI
    {\seq \emp
      {\lam x {A×(B+C)} {\match {\pr2 x}
          b {\inl {\pair {\pr1 x} b}}
          c {\inr {\pair {\pr1 x} c}}}}
      {A×(B+C) → A×B + A×C}}
    {\sumE
      {\seq {x : A×(B+C)}
        {\match {\pr2 x}
          b {\inl {\pair {\pr1 x} b}}
          c {\inr {\pair {\pr1 x} c}}}
        {A×B + A×C}}
      {\prodER
        {\seq {x : A×(B+C)} {\pr2 x} {B+C}}
        {\ax {\seq {x : A×(B+C)} {x} {A×(B+C)}}}
      }
      {\sumIL
        {\seq {Γ ≔ x : A×(B+C), b : B} {\inl {\pair {\pr1 x} b}} {A×B + A×C}}
        {\prodI
          {\seq {Γ} {\pair {\pr1 x} b} {A×B}}
          {\prodEL
            {\seq Γ {\pr1 x} A}
            {\ax {\seq Γ x {A×(B+C)}}}
          }
          {\ax {\seq Γ b B}}
        }
      }
      {\sumIR
        {\seq {Δ ≔ x : A×(B+C), c : C} {\inr {\pair {\pr1 x} C}} {A×B + A×C}}
        {\prodI
          {\seq {Δ} {\pair {\pr1 x} c} {A×C}}
          {\prodEL
            {\seq Δ {\pr1 x} A}
            {\ax {\seq Δ x {A×(B+C)}}}
          }
          {\ax {\seq Γ c C}}
        }
      }
    }
  \end{mathpar}
\end{solution}


\end{document}


%%% Local Variables:
%%% mode: latex
%%% TeX-master: t
%%% End:
